\section{A barrier for Hermite multiplicity retuning}
\label{sec:taper-barrier-audit}
This lemma concerns the normalization loss certificate in the paper's
Cauchy--Hermite construction. It is not a lower bound for SAT, tensor
contraction, or optimal Hermite recovery.

\begin{proposition}
Let $a\ge u\ge1$ be integers, $M>0$, $L=2^a$, and let
$q_0,\ldots,q_{u-1}$ be nonnegative integers.
Put $q_i=0$ for $u\le i\le a$. Suppose
\[
 L\sum_{i=0}^{u-1}2^{-i}q_i>M(2L-1).
\]
For $0\le i<u$, define
\[
 \ell_i=\sum_{j=1}^{a-i}\frac{j}{2^j}(q_i-q_{i+j}).
\]
Then some $i<u$ with $q_i>0$ satisfies
\[
 \ell_i>\frac{M}{4u}.
\]
Consequently, for any valuation-symmetric multiplicity profile satisfying
the degree-recovery condition, some retained node has normalization loss
$-v_2(E_t(0))=D_t\ell_i>D_tM/(4u)$.
\end{proposition}
\begin{proof}
Write $Q=\max_{0\le i<u}q_i$. Since $\sum_{i=0}^{u-1}2^{-i}<2$,
the degree condition gives $Q>M(1-1/(2L))\ge M/2$.
Let
\[
 W_i=\sum_{j=1}^{a-i}j2^{-j}\ge\frac12,\qquad
 d_i=\ell_i/W_i.
\]
Consider the following finite increasing Markov chain. From a state
$i<u$, jump to $i+j$, for $1\le j\le a-i$, with probability
$j2^{-j}/W_i$. Stop at the first state at least $u$. Every jump is
positive, so the stopping time $\tau$ is at most $u$, and
\[
 q_i-\mathbb E[q_{X_{t+1}}\mid X_t=i]=d_i.
\]
Start at a state with $q_i=Q$. Telescoping conditional expectations,
using $q_{X_\tau}=0$, yields
\[
 Q=\mathbb E\sum_{t=0}^{\tau-1}d_{X_t}
 \le u\max_{i<u}d_i.
\]
At a maximizing state, $d_i\ge Q/u>0$ and therefore
$\ell_i=W_id_i\ge Q/(2u)>M/(4u)$. If $q_i=0$, nonnegativity
of all later $q$ would give $\ell_i\le0$, so this state is retained.
The final assertion follows from the exact valuation identity in
the proof of Lemma D.5.
\end{proof}

The proposition does not require a monotone taper. It shows that replacing
the linear taper by another valuation-symmetric profile cannot improve the
\emph{worst normalization-loss bound} from order $MD/u$ to $o(MD/u)$.
It does not rule out a different recovery identity that exploits
cancellation among final Hermite coefficients, an asymmetric selection of
nodes, or a different deformation.

\section{Newton modes: exact formula, but no better logarithmic power}
For a crossing tile, put $h=2^{i+1}$ and
$A=T-y_0+x$, so $v_2(A)\le i$. For
$g(j)=(A-hj)^{-1}$, the forward-difference identity is
\[
 \Delta^b g(0)=
 \frac{b!\,h^b}{\prod_{r=0}^{b}(A-hr)}.
\]
It follows by induction from the definition of $\Delta$, or directly by
partial fractions. On a tile with integer indices $j=0,\ldots,d-1$,
Newton interpolation therefore represents the kernel exactly as
\[
 g(j)=\sum_{b=0}^{d-1}\binom jb\,\Delta^b g(0).
\]
The summands with $b>j$ vanish. This finite expression can replace the
geometric modes $j^b$ by the integral modes $\binom jb$.

Write $v=v_2(A)\le i$. Every denominator factor $A-hr$ has valuation
$v$, since $v_2(h)>v$. Thus
\[
 v_2(\Delta^b g(0))
 =-v+b(i+1-v)+v_2(b!).
\]
Multiplying by the Cauchy prefactor gives the certified weight
\[
 a+1-v+b(i+1-v)+v_2(b!).
\]
In the worst crossing valuation $v=i$, this becomes
\[
 a+1-i+b+v_2(b!)=a+1-i+2b-s_2(b),
\]
where $s_2(b)$ is the number of ones in the binary expansion of $b$.
Hence the weight remains linear in $b$. The factorial contributes only
a constant-factor improvement to the slope, not a superlinear weight.
The finite cap $b<d=D/2$ also does not remove modes at the partition
saddle scale $b\asymp D/\sqrt u$.

This is an exact alternative kernel representation worth investigating
for implementation constants. It does not prove an exponent better than
$O(M/\sqrt u)$, and therefore does not meet the stronger-than-Appendix-D
target.

